\documentclass[a4paper,11pt]{article}

\usepackage{revy}
\usepackage[utf8]{inputenc}
\usepackage[T1]{fontenc}
\usepackage[danish]{babel}
\usepackage{amsmath,amssymb,amsfonts,amsthm,textcomp}


\revyname{MatematikRevyen}
\revyyear{2026}
\version{0.1}
\eta{$2$ minutter}
\status{Færdig}

\title{Dørvagterne 2}
\author{Carl '21}

\begin{document}
\maketitle

\begin{roles}
\role{S1}[Skuespiller] Studerende
\role{S3}[Skuespiller] Studerende
\role{S4}[Skuespiller] Studerende
\role{D1}[Skuespiller] Dørvagt
\role{T}[Taler] Speaker-stemme
\end{roles}

\begin{props}
\item Nogle studiekort
\end{props}


\begin{sketch}
\scene{Samme setting som dørvagterne 1. S1 kommer ind på scenen og går hen til den første. De har denne gang taget flere studiekort med}

\says{S1} Her! Så kan I tjekke dem samtidig

\scene{S1 deler studiekort ud til de først 3-4 vagter. D1 giver det hurtigt tilbage. Ellers går processen for S1 stadig langsomt. S3 kommer ind på scenen og går hen til D1}

\says{D1} Studikort tak

\says{S3} Hvis jeg ikke kan komme ind, så er det fordi mit studiekort ikke virker, sandt?

\says{D1} Joh, det lyder rigtig nok

\says{S3}[viser kortet] Men se her

\says{D1}[tjekker kortet] Ja, den er god nok. Du går bare ind

\scene{S3 går forbi hele rækken. S1 bliver irriteret, smider kortene på jorden og giver op igen}

\scene{lys halvt ned}

\says{T} Stop med de direkte beviser. Gør som S3 her og lær at kontraponer

\scene{lys ned}

\scene{Folk begynder langsomt at forlade scenen}

\scene{lys op igen}

\scene{S4 kommer ind og har lidt travlt. Vagterne er så småt på vej væk, men S4 fanger D1}

\says{S4} Hey! Luk mig lige ind

\says{D1} Hvorfor skulle jeg det?

\says{S4}[tænker] Fordi.. fordi.. hvad hvis du ikke ville lukke mig ind?

\says{D1} Hvad med det?

\says{S4} Det ville være rimelig fucked up, ville det ikke?

\says{D1} Hvorfor det?

\says{S4} Øhhh, så kunne jeg ikke studere?

\says{D1} Ja, og?

\says{S4} Såøhh, ahh, æhh...

\scene{lys halvt ned}

\says{T} Bare lad vær med at spild din tid, med at lede efter den modstrid.

\scene{lys ned}


\end{sketch}
\end{document}
